% Generated draft: requires mathematical and source review.
\documentclass[11pt]{article}
\usepackage[margin=25mm]{geometry}
\usepackage{fontspec}
\setmainfont{Latin Modern Roman}
\usepackage{amsmath,amssymb,mathtools,unicode-math}
\setmathfont{Latin Modern Math}
\usepackage{enumitem,needspace,xurl,hyperref}
\setlength{\parindent}{0pt}
\setlength{\parskip}{5pt}
\setlength{\emergencystretch}{3em}
\newcommand{\N}{\mathbb{N}}
\newcommand{\Z}{\mathbb{Z}}
\newcommand{\Q}{\mathbb{Q}}
\newcommand{\R}{\mathbb{R}}
\newcommand{\C}{\mathbb{C}}
\newcommand{\F}{\mathbb{F}}
\newcommand{\A}{\mathbb{A}}
\newcommand{\PP}{\mathbb{P}}
\newcommand{\E}{\mathbb{E}}
\newcommand{\eps}{\varepsilon}
\newcommand{\dd}{\,\mathrm{d}}
\begin{document}
\section*{UnsolvedMath ID: 30000566}
\textbf{Source problem code:} catalogue entry
\subsection{Definitions and mathematical statement}
Does there exist a continuous path of hypercyclic operators with common hypercyclic vectors?

For a continuous linear operator $T$ on a topological vector space $X$, define
\[
 \operatorname{HC}(T)=
 \{x\in X:\overline{\{T^nx:n\in\N_0\}}=X\}.
\]
The operator is hypercyclic if this set is nonempty. A common hypercyclic vector for a family $(T_t)_{t\in I}$ belongs to
\[
 \bigcap_{t\in I}\operatorname{HC}(T_t).
\]
Thus the intended existence condition is a continuous map $t\mapsto T_t$ into a specified operator space, with every $T_t$ hypercyclic and that intersection nonempty. For a fixed common vector, the iterate used to meet a given open set may depend on $t$ and the open set; no uniform iterate is required.

The source gives no space $X$, scalar field, operator topology, parameter interval, or endpoint conditions. Norm continuity, strong-operator continuity, and other topologies give different questions. It also does not say whether a path must be nonconstant or injective. Allowing a constant path would not express a nontrivial family problem, so nontriviality is an essential missing qualification; it is not silently added as an explicit source hypothesis. Nor is the question replaced by one requiring a dense or residual common-vector set, or a common hypercyclic linear subspace.

Accordingly, the extracted sentence is retained as an existence question requiring recovery of its ambient operator class and path restrictions. The supplied verification calls it internally complete, but these omissions prevent a uniquely determined nontrivial formulation.
\subsection{Short English statement}
The question seeks a continuous operator family sharing at least one vector with dense orbit for every member. The operator topology, underlying space, and a nonconstant-path requirement are missing, so the intended nontrivial path problem is not yet precise.
\subsection{Sources}
\begin{itemize}
\item Local UnsolvedMath registry, permanent ID 30000566, code catalogue entry, ``Paths of Hypercyclic Operators with Common Vectors.'' Proposed year 2006; proposer unspecified.
\item \emph{Mini-Workshop: Hypercyclicity and Linear Chaos} (2007), Oberwolfach Reports, pp.\ 2227--2276, DOI 10.4171/owr/2006/37; \url{https://doi.org/10.4171/owr/2006/37}. The displayed report year 2007 differs from the DOI and proposed year 2006. This is the sole supplied literature reference.
\item Statement status exact; the supplied verification says the statement is clear and internally complete and was retained byte-for-byte. Its missing topology and nontriviality conventions still require source review.
\item Literature triage checked 2026-08-21: open. It generically calls the item an open problem/conjecture/direction and reports no authoritative resolution in the supplied targeted check. Created 2026-08-25; updated 2026-09-16. Current status and precise source scope require independent verification.
\item Registry curation by Ulam AI / UnsolvedMath Contributors, CC BY 4.0; the underlying report retains its supplied provenance. Unreviewed statement-only draft.
\end{itemize}
\end{document}
